\section*{Bab \ref{ch:g12:rlc-oscillations} --- Osilasi Listrik Bebas: RLC}

\begin{solution}{exo:g12:rlc-oscillations:1}
(a) $\sqrt{LC} = \qty{1.0e-4}{s}$: $T_0 = \qty{0.63}{ms}$,
$f_0 \approx \qty{1.6}{kHz}$. (b) $\sqrt{LC} = \qty{3.16e-7}{s}$:
$T_0 = \qty{2.0}{\micro s}$, $f_0 \approx \qty{0.50}{MHz}$ --- berselisih
faktor $\sqrt{10^5} \approx 316$.
\end{solution}

\begin{solution}{exo:g12:rlc-oscillations:2}
$LC$: $(\unit{V.s/A})(\unit{C/V}) = \unit{s} \times \unit{C/A} =
\unit{s} \times \unit{s} = \unit{s^2}$, sehingga $\sqrt{LC}$ bersatuan sekon.
$2\pi$ adalah bilangan murni --- tanpa dimensi untuk diubah.
\end{solution}

\begin{solution}{exo:g12:rlc-oscillations:3}
$Q_0 = CU_0 = \qty{2.8}{\micro C}$; $E = \tfrac12 CU_0^2 =
\tfrac12 \times \num{4.7e-7} \times 36 \approx \qty{8.5}{\micro J}$.
Muatannya berosilasi pada $f_0 = 1/(2\pi\sqrt{LC})$, arusnya terpaut seperempat
\omterm{def:g12:oscillators-and-time:oscillator}{periode} \omterm{def:g12:oscillators-and-time:phase}{fasenya}, \omterm{def:g10:energy-conservation:energy}{energinya} berpindah antara kepingnya dan \omterm{def:g11:magnetic-fields:solenoid}{kumparannya}, dengan
jumlahnya tetap.
\end{solution}

\begin{solution}{exo:g12:rlc-oscillations:4}
$Q_0 = \qty{2.0}{\micro C}$; $f_0 = 1/T_0 = \qty{250}{Hz}$;
$q(T_0/2) = \qty{-2.0}{\micro C}$; $i(0) = 0$ (kosinusnya mendatar pada $t=0$);
$I_0 = 2\pi \times \num{2.0e-6}/\num{4.0e-3} \approx \qty{3.1}{mA}$.
\end{solution}

\begin{solution}{exo:g12:rlc-oscillations:5}
Seluruhnya di \omterm{def:g12:rc-rl-circuits:capacitor}{kapasitornya} ketika $q = \pm Q_0$, $i = 0$ ($t = 0$, $T_0/2$,
$T_0$, \dots): ujung ayunan bandulnya. Seluruhnya di \omterm{def:g11:magnetic-fields:solenoid}{kumparannya} ketika $q = 0$,
$i = \pm I_0$ ($t = T_0/4$, $3T_0/4$, \dots): titik terendah dan tercepat
bandulnya.
\end{solution}

\begin{solution}{exo:g12:rlc-oscillations:6}
$d^2q/dt^2 = -(2\pi/T_0)^2\,Q_0\cos(2\pi t/T_0) = -(2\pi/T_0)^2 q$,
sehingga $L\,d^2q/dt^2 + q/C = q\,[\,1/C - L(2\pi/T_0)^2\,] = 0$ untuk semua $t$
tepat ketika $(2\pi/T_0)^2 = 1/(LC)$, yaitu $T_0 = 2\pi\sqrt{LC}$.
\end{solution}

\begin{solution}{exo:g12:rlc-oscillations:7}
$L = \dfrac{1}{4\pi^2 f_0^2 C} = \dfrac{1}{39.5 \times
(\num{1.98e5})^2 \times \num{2.2e-10}} \approx \qty{2.9}{mH}$.
\end{solution}

\begin{solution}{exo:g12:rlc-oscillations:8}
$E = \tfrac12 CU_0^2 = \qty{5.0e-5}{J}$; $\tfrac12 LI_0^2 = E$ memberikan
$I_0 = \sqrt{2E/L} = \sqrt{\num{1.0e-2}} = \qty{0.10}{A}$, yang pertama kali
tercapai pada $T_0/4 = \tfrac14 \times \qty{0.63}{ms} \approx
\qty{0.16}{ms}$.
\end{solution}

\begin{solution}{exo:g12:rlc-oscillations:9}
$E_C \propto \cos^2(2\pi t/T_0)$, dan $\cos^2$ berulang tiap
\emph{setengah} putaran: dua kali pengisian \omterm{def:g12:rc-rl-circuits:capacitor}{kapasitornya} tiap \omterm{def:g12:oscillators-and-time:oscillator}{periode}, sehingga
$2f_0$. Pada $t = T_0/8$ \omterm{def:g12:oscillators-and-time:phase}{fasenya} $\pi/4$:
$\cos^2(\pi/4) = \tfrac12$, sehingga $E_C = E_L = E/2$.
\end{solution}

\begin{solution}{exo:g12:rlc-oscillations:10}
Nisbahnya $3.0/4.0 = 0.75$ tiap \omterm{def:g12:rlc-oscillations:regimes}{pseudoperiode}: puncak keenamnya
$4.0 \times 0.75^6 \approx \qty{0.71}{V}$. \omterm{def:g10:energy-conservation:energy}{Energinya} $\propto$ kuadrat
\omterm{def:g12:oscillators-and-time:oscillator}{amplitudonya}: $0.75^2 \approx 56\%$ bertahan tiap \omterm{def:g12:oscillators-and-time:oscillator}{periode}.
\omterm{def:g12:rlc-oscillations:regimes}{Pseudoperiodenya} tetap $\approx T_0 = 2\pi\sqrt{LC}$: \omterm{def:g12:oscillators-and-time:damping}{peredaman} ringan
menyusutkan \omterm{def:g12:oscillators-and-time:oscillator}{amplitudonya}, hampir tidak iramanya.
\end{solution}

\begin{solution}{exo:g12:rlc-oscillations:11}
$T_0 = 2\pi\sqrt{0.50 \times \num{2.0e-5}} = 2\pi \times
\num{3.16e-3} \approx \qty{20}{ms}$. Kamusnya: $k = m/(LC) =
0.50/\num{1.0e-5} = \qty{5.0e4}{N/m}$; periksalah
$2\pi\sqrt{m/k} = 2\pi\sqrt{\num{1.0e-5}}$ --- \qty{20}{ms} yang sama.
\end{solution}

\begin{solution}{exo:g12:rlc-oscillations:12}
$C = \qty{500}{pF}$: $f_0 = 1/(2\pi\sqrt{\num{1.0e-13}}) \approx
\qty{0.50}{MHz}$; $C = \qty{50}{pF}$: $f_0 \approx \qty{1.6}{MHz}$.
Pada $L$ yang tetap, $f_0 = (2\pi\sqrt{L})^{-1} C^{-1/2} \propto 1/\sqrt{C}$:
$C$ sepuluh kali lipat, hanya $\sqrt{10} \approx 3.2$ pada $f_0$ --- seperti yang ditemukan.
\end{solution}

\begin{solution}{exo:g12:rlc-oscillations:13}
$0.95^n = \tfrac12$: $n = \ln 2/\lvert\ln 0.95\rvert \approx 13.5$,
jadi sekitar $14$ \omterm{def:g12:oscillators-and-time:oscillator}{periode}. \omterm{def:g10:signals-and-waves:amplitude}{Amplitudonya} $\propto \sqrt{E}$: turun sampai
$1/\sqrt{2} \approx 0.71$ nilai awalnya. Pada \qty{1.0}{MHz},
$T_0 = \qty{1.0}{\micro s}$: sekitar \qty{14}{\micro s}.
\end{solution}

\begin{solution}{exo:g12:rlc-oscillations:14}
Susutnya tiap daur $0.020 \times \qty{1.0e-9}{J} = \qty{2.0e-11}{J}$, pada
\num{32768} daur tiap sekon: $P = \num{32768} \times \num{2.0e-11}
\approx \qty{6.6e-7}{W}$ --- di bawah semikrowatt, bertahun-tahun dengan sebuah
sel kancing. Harus selaras, sebab hanya dorongan pada \omterm{def:g12:oscillators-and-time:phase}{fase} yang tepat yang
\emph{menambahkan} \omterm{def:g10:energy-conservation:energy}{energi} (\omterm{def:g12:oscillators-and-time:resonance}{resonansi}); jika tidak selaras ia akan mengerem ayunannya.
\end{solution}

\begin{solution}{exo:g12:rlc-oscillations:15}
$C = \dfrac{1}{4\pi^2 f_0^2 L} = \dfrac{1}{39.5 \times
(\num{2.5e7})^2 \times \num{1.0e-6}} \approx \qty{41}{pF}$.
Resistornya seharusnya mendorong simpalnya ke \omterm{def:g12:rlc-oscillations:regimes}{rezim aperiodik}: muatannya
lalu mengendap tanpa melampaui --- tanpa osilasi, tanpa dentingan.
\end{solution}

\begin{solution}{pb:g12:rlc-oscillations:1}
\textbf{1.} Hukum simpalnya: $u_L + u_C = 0$, dengan $u_C = q/C$,
$u_L = L\,di/dt$, $i = dq/dt$: $L\,d^2q/dt^2 + q/C = 0$.

\textbf{2.} $d^2q/dt^2 = -(2\pi/T_0)^2 q$, sehingga persamaannya berbunyi
$q\,[\,1/C - L(2\pi/T_0)^2\,] = 0$, yang benar untuk semua $t$ tepat ketika
$T_0 = 2\pi\sqrt{LC}$.

\textbf{3.} $LC$: $(\unit{V.s/A})(\unit{C/V}) = \unit{s} \times
\unit{C/A} = \unit{s^2}$ --- $\sqrt{LC}$ adalah waktu.

\textbf{4.} $LC = \qty{6.6e-14}{s^2}$: $T_0 = 2\pi \times \num{2.57e-7}
\approx \qty{1.6}{\micro s}$, $f_0 \approx \qty{620}{kHz}$ --- di dalam
\qtyrange{530}{1700}{kHz}.

\textbf{5.} $C = \dfrac{1}{4\pi^2 f_0^2 L} = \dfrac{1}{39.5 \times
\num{1.0e12} \times \num{2.0e-4}} \approx \qty{127}{pF}$.

\textbf{6.} Rumus yang sama pada \qty{530}{kHz}: $C \approx \qty{451}{pF}$.

\textbf{7.} Pada \qty{1700}{kHz}: $C \approx \qty{44}{pF}$.

\textbf{8.} $C = \qty{480}{pF}$: $f_0 \approx \qty{514}{kHz}$;
$C = \qty{20}{pF}$: $f_0 \approx \qty{2.5}{MHz}$ --- seluruh pitanya,
dengan sisa di kedua ujungnya.

\textbf{9.} $f_0 = (2\pi\sqrt{L})^{-1}C^{-1/2} \propto 1/\sqrt{C}$;
menggandakan $f_0$ menuntut membagi $C$ dengan $4$.

\textbf{10.} Karena $f_0 \propto 1/\sqrt{C}$, langkah $C$ yang sama menggeser
$f_0$ perlahan pada $C$ yang besar dan cepat pada $C$ yang kecil: jarak tertentu
antarstasiun lalu hanya menuntut sedikit putaran tombolnya --- stasiunnya berdesakan di
ujung yang berfrekuensi tinggi ($C$ rendah).

\textbf{11.} $Q_0 = CU_0 = \num{1.27e-10} \times 0.020 \approx
\qty{2.5e-12}{C}$; $E = \tfrac12 CU_0^2 \approx \qty{2.5e-14}{J}$.

\textbf{12.} $\tfrac12 LI_0^2 = E$: $I_0 = U_0\sqrt{C/L} =
0.020\sqrt{\num{6.35e-7}} \approx \qty{16}{\micro A}$.

\textbf{13.} $0.96^n = \tfrac12$: $n = \ln 2/\lvert\ln 0.96\rvert
\approx 17$ daur.

\textbf{14.} Pada \qty{1.00}{MHz}, $T_0 = \qty{1.0}{\micro s}$: dentingannya
menyusut separuh dalam sekitar \qty{17}{\micro s}. Satu dentingan mati dalam beberapa mikrosekon;
musiknya menuntut simpalnya disuapi terus-menerus.

\textbf{15.} \omterm{def:g10:signals-and-waves:wave}{Gelombangnya} mendorong jutaan kali tiap sekon; hanya ketika
\omterm{def:g12:oscillators-and-time:oscillator}{frekuensinya} menyamai $f_0$ dorongannya tiba selaras lalu membangun
ayunannya --- \omterm{def:g12:oscillators-and-time:resonance}{resonansi}; stasiun yang tidak tertala mendorong sama seringnya melawan maupun
searah dengannya. Di pemancarnya, sebuah penguat mengembalikan \omterm{def:g10:energy-conservation:energy}{energi} yang terlesap
tiap daur: sebuah \omterm{def:g12:oscillators-and-time:oscillator}{osilator} terpelihara.

\textbf{16.} $LC = \qty{2.54e-14}{s^2}$: $k = m/(LC) =
0.20/\num{2.54e-14} \approx \qty{7.9e12}{N/m}$.

\textbf{17.} Sekitar $10^8$ kali pegas mobilnya: tidak ada sistem
pegas--massa di atas meja yang mencapai megahertz --- \omterm{def:g12:oscillators-and-time:spring}{kekakuan} dan massa
sehari-hari terlalu lamban.

\textbf{18.} $m = \dfrac{k}{(2\pi f_0)^2} = \dfrac{100}
{(\num{6.28e6})^2} \approx \qty{2.5e-12}{kg}$ --- beberapa nanogram, sebutir
debu.

\textbf{19.} Sebutir hablur kuarsa: sekeping seringan bulu dan amat kaku
yang bernyanyi pada \omterm{def:g12:rlc-oscillations:period}{frekuensi alaminya} --- penyimpan waktu yang ditemui lebih awal
tahun ini.

\textbf{20.} Dengan $L = \qty{0.20}{mH}$ dan \qtyrange{20}{480}{pF}
simpalnya menyapu sekitar \qtyrange{0.51}{2.5}{MHz}, meliputi pitanya;
$C \approx \qty{127}{pF}$ memilih \qty{1.00}{MHz}; jika dibiarkan sendiri dentingannya
menyusut separuh dalam $\approx \qty{17}{\micro s}$, tetapi \omterm{def:g10:signals-and-waves:wave}{gelombang} milik stasiunnya,
yang beresonansi selaras, menyuapi kembali simpalnya --- wartanya bermain dengan \omterm{def:g10:energy-conservation:energy}{energi}
\omterm{def:g10:signals-and-waves:wave}{gelombangnya} sendiri.
\end{solution}
